% !PRINTED 492
\weberpage{492}

Dann k\"onnen die Producte $\omega \omega_1, \omega \omega_2, \ldots, \omega \omega_m$ mit H\"ulfe der Gleichungen (2) in die Form gesetzt werden
\[
\omega \omega_r = c_{r,1} \omega_1 + c_{r,2} \omega_2 + \cdots + c_{r,m} \omega_m \quad r = 1, 2, \ldots, m,
\]
worin die Co\"efficienten $c_{s,r}$ ganze rationale Zahlen sind. Wenn man aus diesen Gleichungen aber die $\omega_1, \omega_2, \ldots, \omega_m$ eliminirt, so folgt
\[
\begin{vmatrix}
c_{1,1} - \omega, & c_{1,2}, & \ldots, & c_{1,m} \\
c_{2,1}, & c_{2,2} - \omega, & \ldots, & c_{2,m} \\
\cdots & \cdots & \cdots & \cdots \\
c_{m,1}, & c_{m,2}, & \ldots, & c_{m,m} - \omega
\end{vmatrix}
= 0,
\]
was entwickelt die Form enth\"alt
\[
\omega^m + C_1 \omega^{m-1} + \cdots + C_m = 0,
\]
worin die $C_1, C_2, \ldots, C_m$ gleichfalls ganze rationale Zahlen sind. Dies aber zeigt, dass $\omega$ eine ganze Zahl ist, wie bewiesen werden sollte.

3. \gesperrt{Ist $f(x)$ eine im K\"orper der rationalen Zahlen irreducible Function, und ist eine Wurzel $\Theta$ von $f(x) = 0$ eine ganze Zahl, so sind alle Wurzeln von $f(x)$ ganze Zahlen.}

Denn wenn eine rationale Function $F(x)$ f\"ur $x = \alpha$ verschwindet, so ist $F(x)$ durch $f(x)$ theilbar, und alle Wurzeln von $f(x)$ sind zugleich Wurzeln von $F(x)$ (Bd.\,I, \S.\,141). Wenn nun $\alpha$ eine ganze Zahl ist, so giebt es eine Function
\[
F(x) = x^m + A_1 x^{m-1} + \cdots + A_m
\]
mit ganzzahligen Co\"efficienten $A_1, \ldots, A_m$, die f\"ur $x = \alpha$ verschwindet, und $F(x)$ verschwindet also auch f\"ur alle anderen Wurzeln von $f(x)$, die sonach alle ganze Zahlen sind.

Ist
\[
f(x) = x^n + a_1 x^{n-1} + \cdots + a_n,
\]
so sind die $a_1, a_2, \ldots, a_n$ durch Multiplication und Addition aus den Wurzeln von $f(x)$ zusammengesetzt und sind also nach 2. ganze rationale Zahlen. Daraus folgt:

4. \gesperrt{Ist $\Theta$ eine ganze algebraische Zahl, so hat die Gleichung niedrigsten Grades $f(\Theta) = 0$ [\S.\,132, (2)] ganzzahlige Co\"efficienten.}
